\documentclass[10pt,a4paper]{amsart}
\usepackage[margin=19mm]{geometry}
\usepackage{amsmath,amssymb,mathtools}
\usepackage[scr=rsfs,cal=euler]{mathalfa}
\usepackage{xfrac}
\usepackage[unicode,hidelinks,bookmarks=false]{hyperref}
\newcommand{\ie}{i.e.\ }
\newcommand{\cf}{cf.\ }
\newcommand{\resp}{respectively\ }
\newcommand{\Subsection}{\subsection}
\providecommand{\nobreakdash}{\nobreak\hbox{-}}
\setlength{\emergencystretch}{2em}
\multlinegap0pt
\usepackage{amssymb}
\usepackage{float}
\newtheorem{theorem}{Theorem}
\title{List and total colorings of multiset permutation graphs}
\author{Italo J. Dejter}
\date{Source: arXiv:2602.09965v1 (CC BY 4.0)}
\begin{document}
\maketitle

\noindent\emph{Source and licence.} List and total colorings of multiset permutation graphs. Version arXiv:2602.09965v1 (CC BY 4.0), distributed by arXiv under the Creative Commons Attribution 4.0 International licence, http://creativecommons.org/licenses/by/4.0/, as recorded in the arXiv licence metadata for this exact version and shown on the abstract page https://arxiv.org/abs/2602.09965v1. Full licence notice: \texttt{licenses/arxiv-2602.09965-CC-BY-4.0.txt}.\par\smallskip

\noindent\textit{Editorial scope.} Counted unit: the article's theorem choo (source0.tex lines 232--235). The article's complete original proof of it is delivered with it (source0.tex lines 237--242, one \texttt{proof} environment of 6 lines). No mathematics is written, repaired or re-derived: the statement and the proof are the article's own lines, in the article's own notation, and no retained line is substituted or re-flowed.  Retained uncounted as the source-local context the proof needs: lines 214--228.  Nothing was omitted from any retained span, so the package carries no omission record.  No retained line contains a citation, so the wrapper carries no bibliography and no bibliography entry.  No retained line contains a figure environment, an image-inclusion command or drawing code, so no image is delivered and the package declares no asset dependency.  Editorial formatting only, in addition: an AMS \texttt{amsart} preamble that restates the article's own declarations and macros used by the retained lines, namely the environments \texttt{theorem} on separate counters, together with the standard packages those lines need.  The retained environments print in this document as Theorem 1 in order of appearance, whereas the article prints the same items with the numbering of its own full document.  The complete native source of the article as distributed in the arXiv e-print for this exact version is bundled unchanged as \texttt{sources/194361/source0.tex}.
\begin{proof}
A proper edge coloring of $ST_k^\ell$ is given as follows. Let $i\in[k\ell]\setminus\{0\}$. Then, 
each edge 
\begin{eqnarray}\label{notice}
(v,w)=(v_0v_1\cdots v_i\cdots v_{k\ell-1},\;v_iv_1\cdots v_0\cdots v_{k\ell-1}),
\mbox{ where }v_i\ne v_0,
\mbox{ is assigned color }i.
\end{eqnarray} 
The resulting edge color assignment also is to include that the edge
\begin{eqnarray}\label{not2}\begin{array}{ll}
(v,w)=(v_0v_1\cdots v_{k\ell-1},\;v_1v_0\cdots v_{k\ell-1}),\mbox{ where }v_1\ne v_0,&\mbox{ is assigned color }1,\mbox{ and}\\
(v,w)=(v_0v_1\cdots v_{k\ell-1},\;v_{k\ell-1}v_1\cdots v_{0}),\mbox{ where }v_{2k-1}\ne v_0,&\mbox{ is assigned color }k\ell-1.
\end{array}\end{eqnarray} 
As a result, the edges incident to each vertex of $ST_k^\ell$ have pairwise different colors, so that $ST_k^\ell$ has a proper edge coloring.
\end{proof}

\begin{theorem}\label{choo}
Let $k>1$ and let $\ell>2$. Then, the graph $ST_k^\ell$ is $(\ell-1)$-choosable. However, no $(\ell-1)$-list coloring of $ST_k^\ell$ contains only efficient color classes.
%, meaning that some class of each $(\ell-1)$-list coloring of $ST_k^\ell$ is not an E-set.
\end{theorem}

\begin{proof}
There are $\ell-1$ occurrences of the first entry $v_0$ in any vertex $v=v_0\cdots v_{k\ell-1}$ of $ST_k^\ell$. The remaining $(k-1)\ell$ entries in $\{1,\ldots,k\ell-1\}$ determine the edges incident to $v$; these edges are colored according to displays (\ref{notice}) and (\ref{not2}). The said $\ell-1$ occurrences constitute a list $L(v)$, $\forall v\in V(ST_k^\ell)$. Adjacent vertices $v=v_0\cdots v_{k\ell-1}$ and $w=w_0\cdots w_{k\ell-1}$ have $L(v)\cap L(w)=\emptyset$, since $v_0\neq w_0$, where $w_0=v_j$ for some $j\in\{1,\dots,k\ell-1\}$ and $v_j\neq v_0$, that is: $v_j$ must have a different value in $\{0,\ldots,k-1\}$ than that of $v_0$. Thus, $ST_k^\ell$ is $(\ell-1)$-choosable.

To show that no $(\ell-1)$-list coloring of $ST_k^\ell$ contains only efficient color classes, note that each vertex $w$ at distance 2 from say vertex $v=00\cdots 0v_\ell\cdots,v_{k\ell-1}$
starts with $w_0=0$ and only has one entry in $\{w_1,\ldots,w_{\ell-1}\}$ distinct from 0. But there are $(\ell-1)(\ell-2)$ such vertices $w$ and only $\ell-1$ colors availabe, so by the Pigeonhole Principle, the last sentence of the statement holds in this case. A similar argument holds by setting $v_0\ne 0$ and taking any permutation of the remaining entries of $v$.
\end{proof}

\end{document}
